%*******************************************************
% Resume / Abstract
%*******************************************************
\pdfbookmark[1]{Résumé}{Resume}
\begingroup
\let\clearpage\relax
\let\cleardoublepage\relax

\chapter*{Résumé}

% ============================ TODO texte (~250 mots) ===============================
% Le contenu classicthesis d'origine (« Short summary of the contents in English,
% a great guide by Kent Beck... ») a ete retire: il etait encore en place et
% partait a l'impression.
%
% Huit elements a faire tenir en 250 mots, dans cet ordre:
%
%  1. LA MATRICE DE COVARIANCE comme fil unique du dossier: le seul objet que
%     les trois chapitres de contribution estiment, contraignent ou apprennent.
%
%  2. LES DEUX FORCES qui s'exercent sur elle, et le fait qu'elles ne sont pas
%     independantes: robustesse (donnees non gaussiennes, heterogenes, a
%     queues lourdes) et geometrie (cone SPD, invariance affine).
%
%  3. LES TROIS CAPTEURS, NOMMEMENT --- sonar a croix de Mills, sonar frontal
%     multi-pings, series temporelles d'images RSO. Des instruments concrets se
%     lisent bien mieux que « applications radar et sonar ».
%
%  4. LE SECOND AXE: la geometrie comme contrainte d'architecture en
%     apprentissage profond (la donnee, le batch, les poids).
%
%  5. L'AFFIRMATION METHODOLOGIQUE, a plat: choisir une moyenne, c'est choisir
%     une geometrie, donc un modele statistique.
%
%  6. L'APPAREIL DE REPRODUCTIBILITE: chaque figure de simulation porte son
%     commit, sa graine et ses parametres, et se rejoue. Aucun memoire
%     comparable ne peut l'annoncer: c'est un element differenciant.
%
%  7. UNE PHRASE sur ce que les perspectives ouvrent (attention sur le cone SPD,
%     phase geometrique).
%
%  8. MOTS-CLES en fin de resume si l'ecole doctorale les demande.
% ====================================================================================

\vfill

\chapter*{Abstract}

% ============================ TODO texte (~250 mots) ===============================
% Version anglaise du bloc ci-dessus. NE PAS traduire mot a mot: l'anglais d'un
% resume d'hdr est plus direct que le francais, et la phrase 5 gagne a etre
% plus seche encore.
% ====================================================================================

\vfill

\endgroup

\vfill
